% sections/abstract-uk.tex --- Ukrainian abstract + keywords (9 pt, 1 cm indent).

\begin{otherlanguage}{ukrainian}
\begin{citabstract}
\noindent\textbf{Богусевич О.~О., Свистунов А.~О. PARCS-Agent: міст від агентів великих мовних моделей до розподілених обчислювальних кластерів через протокол MCP.}
Сучасні агенти на основі великих мовних моделей використовують зовнішні інструменти для пошуку інформації, виконання коду та звернень до API, проте кожний виклик інструмента залишається локальною синхронною функцією. Обчислювально-важкі підзадачі --- методи Монте-Карло, перебір гіперпараметрів, комбінаторний пошук --- змушені виконуватися в єдиному потоці агента, і не існує загальноприйнятого механізму отримання додаткових обчислювальних ресурсів під час виконання задачі. У роботі запропоновано архітектурний шаблон PARCS-Agent, у якому сервер протоколу контексту моделі (Model Context Protocol, MCP) надає інтерфейс до Kubernetes-кластера системи PARCS і експонує невеликий набір інструментів, що дозволяють агентові подавати довільну логіку обчислень мовою C\# під час виконання. Платформа компілює подану реалізацію за допомогою Roslyn, розподіляє її між пулом робочих демонів засобами автомасштабування KEDA та передає JSON-результати між послідовними паралельними шарами. Агент не пише паралельного коду --- лише логіку одного робітника, --- що уникає відомих складнощів, які виникають у мовних моделей при роботі з MPI, OpenMP та CUDA. Передача великих вхідних даних здійснюється через сторонній канал спільного сховища з контентно-адресованим кешуванням. Підхід проілюстровано на прикладі обчислення Value-at-Risk методом Монте-Карло; кількісне оцінювання залишене для подальших публікацій.

\medskip
\noindent\textbf{Ключові слова:} ШІ-агенти; великі мовні моделі; протокол контексту моделі;
розподілені обчислення; паралельні обчислення; використання інструментів великими мовними моделями;
Kubernetes; агентна інфраструктура.
\end{citabstract}
\end{otherlanguage}
